%/////////
\begin{frame}[c,allowframebreaks,fragile]{Citations}
%
\begin{block}{\texttt{\textbackslash{}ccite} command\textemdash{}e.g., \ccite{bibtex-patashnik88}}
\begin{verbatim}
\ccite{bibtex-patashnik88}
\end{verbatim}
\begin{itemize}
	\item to be used instead of standard \texttt{\textbackslash{}cite} command
	\item its colour comes from the surface: light on a dark header, as here above
	\item can be used as a note\ccite{bibtex-patashnik88}, with no space before\textemdash{}faint in running text, by design
	\item since style 1.6.1 it takes the optional note of \texttt{\textbackslash{}cite}\ccite[p.~42]{bibtex-patashnik88}
	\begin{itemize}
		\item as here above, from \verb|\ccite[p.~42]{bibtex-patashnik88}|
	\end{itemize}
\end{itemize}
\end{block}
%
\begin{exampleblock}{\texttt{\textbackslash{}ccite} in an \texttt{exampleblock} header\textemdash{}e.g., \ccite[p.~42]{bibtex-patashnik88}}
\begin{itemize}
	\item light, as on every dark header
\end{itemize}
\end{exampleblock}
%
\begin{alertblock}{\texttt{\textbackslash{}cccite} command\textemdash{}e.g., \cccite{bibtex-patashnik88}}
\begin{itemize}
	\item since style 1.1 the same as \texttt{\textbackslash{}ccite}
	\begin{itemize}
		\item a citation marker is already as small as it can be read
		\item kept so that existing decks build as before
	\end{itemize}
\end{itemize}
\end{alertblock}
%
\end{frame}
